\scanpage{378}
\mspage{294}
% archive sheet number: 377

le fait que $\unsure{\alpha|(*)}$ soit iso équivaut à dire
que le foncteur \deleted{d'inclusion} canonique
\[
J\overset{\mathrm{def}}{=}\underline S(\pi)_{/x}\hookrightarrow M_{\pi/x}=I
\tag*{\deleted{$\unsure{(*}$}}
\]
est \uline{cofinal}. (cf 4.4.2 p. 79). \deleted{\unsure{[}}Mais
on voit tout de suite que ce foncteur
est un foncteur d'inclusion d'une sous-catégorie \inserted{str.} pleine [, stable par $\varinjlim$ gdes
devant $\pi'_0$, de cardinal $\leq\pi$].

Notons maintenant que la flèche
$(**)'$ similaire à $(*)$, \unsure{où|avec} $\underline S(\pi)$
remplacée par $\underline S_\alpha(\pi)$, est bel et bien
un iso, puisque $x\in\operatorname{Ob}\underline S\subset\operatorname{Ob}\underline S_\alpha$,
et que $\underline S_\alpha$ est $\pi$-accessible. Cela
s'interprète par le fait que le foncteur
d'inclusion
\[
J_\alpha\overset{\mathrm{def}}{=}\underline S_\alpha(\pi)_{/x}\unsure{\hookrightarrow|\lhook\joinrel\relbar}M_{\pi/x}=I
\]
est cofinal. C'est encore une inclusion
\deleted{\unsure{pl.}} d'une sous-catégorie \inserted{$J_\alpha$} strict\textsuperscript{mt} pleine \inserted{de $I$},
stable par $\varinjlim$ g\textsuperscript{r}des devant $\pi'_0$ et de
cardinal $\leq\pi$, et on a
\[
J=\bigcap J_\alpha,
\]
Donc la conclusion voulue résultera du
